\section{K-033: an exactly evaluated four-dimensional component normalization}

This calculation concerns a component of the actual ten-dimensional
tensor quotient fixed by the originally targeted subgroup
$\mu_{2^{10-1}}=\mu_{512}$.  Its normalization and top Hilbert density
are evaluated completely.  Its possible conductor contributions are
not determined, and no global pole is inferred.

\subsection*{Construction and complete fixed representation}

Take the dimension-nine K-020 tensor and its selected output edge.
The critical mass of that edge is $\rho=293/512$.  Replace it by
the four edges of K-015a, introducing one new coordinate in each
factor.  Denote the resulting critical dimension-ten tensor by $T$.

\begin{lemma}[K-033a]
The full product-linear stabilizer of $T$ is
$S\times\mu_{1024}$, with faithful determinant character.  In
particular its effective special-linear stabilizer is trivial.
The grading subgroup $\mu_{512}$ has a four-dimensional fixed
component $Y$ that is smooth at this quotient point.

For a fixing element $h$ inducing a primitive grading element of
order $512$, each factor spectrum has one double eigenvalue and
eight simple eigenvalues.  The complete fixed representation is
\[
 V_{10}^h=(A_2\otimes B_2\otimes C_2)\oplus A_2^{\oplus2}
                \oplus B_2\oplus C_2^{\oplus2}\oplus\mathbb C^{19}.
                                                        \tag{K-033a}
\]
The five vector masses, in this order, are
\[
 \frac{55}{512},\ \frac{164}{512};\qquad
 \frac{219}{512};\qquad
 \frac{200}{512},\ \frac{19}{512}.                       \tag{K-033b}
\]
The scalar summands retain their outside torus characters.
\end{lemma}

\begin{proof}
The selected output of K-020 has pair-deletion coverage.  K-015a
therefore proves that the doubled tensor has Lie stabilizer $S$.
Its new masses are $\rho/2,\rho/2,\rho/2,1-\rho/2$, with all other
masses unchanged.  They are strictly positive; freeness gives
identity reduced matrices after taking square roots.  The incidence
argument K-006 gives the cyclic diagonal factor of order $1024$
with faithful determinant character, since $\rho$ has odd numerator
and reduced denominator $512$.

For completeness, lift an old primitive generator to modulus $1024$
so that every old edge except the selected edge has weight zero,
and the selected edge has weight $512$.  Such a lift is supplied
by the same K-006 character-lattice calculation and is retained
explicitly in the certificate.  The new weight in each factor is
the negative sum of the other two selected old weights.  Hence
it differs from its selected old weight by $512$.  The four new
edges have weight zero.

Reducing modulo $512$ collapses only the new coordinate onto its
selected old coordinate.  The old spectra were simple, so the
new spectra modulo $1024$ are simple and have exactly one
antipodal pair in each factor.  Every fixed triple collapses to
an old support edge.  An unselected old edge shares at most one
selected coordinate; replacing that coordinate by its new partner
changes the weight by $512$, so gives no further fixed triple.
Over the selected edge, a triple is fixed exactly when an odd
number of its coordinates are new.  These are precisely the four
replacement edges.  Thus the complete fixed support modulo
$1024$ is the constructed support, and the centralizer in the
full stabilizer is its diagonal stabilizer.

Each of the five unselected edges incident to one selected
coordinate gives a nonzero projection onto exactly one of the
three paired subspaces.  K-016 now gives the full stabilizer.
The quotient is consequently smooth at the indicated point.

Modulo $512$, the selected old edge expands to the entire $2^3$
block, and the five other incident edges expand to five vectors.
The other $3\cdot9-8=19$ edges remain scalar.  This proves
(K-033a).  Its dimension is $37$, whereas the fixed
special-linear Lie algebra has dimension
$3(2^2+8-1)=33$.  Exactness of finite-group fixed vectors in the
quotient tangent sequence gives the four-dimensional smooth fixed
germ.  The masses (K-033b) are the corresponding five entries in
the exact K-020 first-step mass vector.
\end{proof}

The proof of K-032 applies without change to this complete fixed
representation: equality of nullcones, finite restriction, and the
unique double eigenvalue prove that
\[
 B=\mathbb C[V_{10}^h]^{C_{SL_{10}^3}(h)}
\]
is the normalization of $\mathbb C[Y]$.  It has only degrees
divisible by $512$.  Put $z=t^{512}$ and
$F(z)=\sum_{a\ge0}\dim B_{512a}z^a$.

\subsection*{The exact Hilbert series}

\begin{theorem}[K-033b]
The normalized component has Hilbert series
\[
 F(z)=\frac{1+91067z+429033z^2+416212z^3+78249z^4}
                  {(1-z)^4(1+z)}.                     \tag{K-033c}
\]
Equivalently, for every integer $a\ge0$,
\[
 \dim B_{512a}
  =\frac{507281}{6}a^3+6410a^2+\frac{338}{3}a
                     +\frac34+\frac{(-1)^a}{4}.        \tag{K-033d}
\]
Its normalized Hilbert density is the three-adic unit $507281$.
In the local parameter $u=1-t^{512}$,
\[
 H_B(t)=\frac{507281}{u^4}-\frac{1001742}{u^3}
       +\frac{1145421}{2u^2}-\frac{312995}{4u}+O(1).
                                                        \tag{K-033e}
\]
In particular every Laurent coefficient of this normalization
series is integral at three.
\end{theorem}

\begin{proof}
Formula K-032f, with $r=2$ and the masses (K-033b), becomes
\[
 \dim B_{512a}=
 \sum_{j=0}^{55a}\sum_{v=0}^{19a}
 g\bigl((512a-j,293a+j),(512a,293a),
                       (512a-v,293a+v)\bigr).          \tag{K-033f}
\]
All three partitions have size $q=805a$.

We use the classical highest-weight generating series for a binary
trilinear form:
\[
 C(t;x,y,z)=
 \frac{1+t^3xyz}
 {(1-txyz)(1-t^2x^2)(1-t^2y^2)(1-t^2z^2)(1-t^4)}.
                                                        \tag{K-033g}
\]
Its coefficient of $t^q x^b y^c z^d$ is the multiplicity of
$\operatorname{Sym}^b\mathbb C^2\otimes
 \operatorname{Sym}^c\mathbb C^2\otimes
 \operatorname{Sym}^d\mathbb C^2$ in
$\operatorname{Sym}^q((\mathbb C^2)^{\otimes3})$.
This is the binary-trilinear covariant series displayed in
Holweck--Luque--Thibon, Section 5.2, printed p.~16.

One may verify (K-033g) independently without assuming a generating
set of covariants.  Define
\[
 \mathcal W_x f(x)=\frac{x f(x)-x^{-1}f(x^{-1})}{x-x^{-1}}.
\]
It maps $x^b$ to the character of $\operatorname{Sym}^b\mathbb C^2$.
Direct multiplication of denominators gives the rational identity
\[
 \mathcal W_x\mathcal W_y\mathcal W_z C(t;x,y,z)
   =\prod_{\epsilon_1,\epsilon_2,\epsilon_3\in\{1,-1\}}
           (1-tx^{\epsilon_1}y^{\epsilon_2}z^{\epsilon_3})^{-1}.
                                                        \tag{K-033h}
\]
The right side is the character series of the symmetric algebra.
Independence of irreducible characters proves the claimed
multiplicity interpretation.  The retained verifier checks this
identity exactly as a rational function, in addition to small
independent symmetric-group character checks.

For the partitions in (K-033f), their three $SL_2$ highest weights
are
\[
 b=219a-2j,\qquad c=219a,\qquad d=219a-2v.
\]
In (K-033g), let $w\ge0$ be the exponent of $txyz$, and let
$\varepsilon\in\{0,1\}$ choose the numerator term.  The three
quadratic generator exponents are forced by
$b=w+\varepsilon+2b'$, and its two analogues.  If $D\ge0$
denotes the exponent of $t^4$, the total degree condition is
\[
 w=-74a-j-v+2D.
\]
Thus $w\equiv j+v\pmod2$, and its nonnegativity already ensures
$D\ge0$.  The quadratic exponents are nonnegative exactly when
\[
 0\le w\le 219a-2\max(j,v)-\varepsilon,
 \qquad \varepsilon\equiv a-j-v\pmod2.
\]
The latter congruence determines $\varepsilon$ uniquely.  Counting
the admissible $w$ gives the following exact two-row coefficient:
\[
 g=\left\lfloor\frac{219a}{2}\right\rfloor-\max(j,v)+1
             -\mathbf1_{\{a\text{ even},\ j+v\text{ odd}\}}.
                                                        \tag{K-033i}
\]
The indicated ranges have $\max(j,v)\le55a$, so no negative-count
truncation is needed.  The case $a=0$ also gives the correct value.

Write $J=55a$ and $V=19a$, so $J\ge V$.  The required elementary
sums are
\[
 \sum_{j=0}^{J}\sum_{v=0}^{V}\max(j,v)
  =\frac{(V+1)J(J+1)}2+\frac{V(V+1)(V+2)}6.
\]
For even $a$, both endpoints are even, and the number of pairs
with $j+v$ odd is $(JV+J+V)/2$.  Substitution in (K-033i) proves
(K-033d), separately for the two parities.  Summing the four
polynomial sequences and the parity sequence gives (K-033c).
Expanding it at $z=1-u$ proves (K-033e).  Its only denominators
are powers of two, and the remaining regular expansion also has
only powers of two as denominators, since the unreduced rational
expression uses only $(1-z)^4(1+z)$.
\end{proof}

\subsection*{The actual component contribution and the remaining gap}

The normalization quotient has dimension at most three.  Therefore
the actual component $Y$ has the same top density $507281$.
Let $E(\xi)$ be its generic moving inverse Euler factor as in
K-032k.  It is a three-adic unit.  The component's order-four local
coefficient is exactly
\[
                            507281 E(\xi).             \tag{K-033j}
\]
K-031 implies that $Y$ lies in a cyclic orbit of three components,
because it has dimension four and is not the chain curve.  The
sum of that orbit's order-four coefficients is
$3\cdot507281 E(\xi)$, of three-adic valuation exactly one.

The lower terms in (K-033e) belong to the normalization Hilbert
series, not automatically to the actual component or its moving
localization class.  In particular we do not infer from (K-033c)
that $Y$ has no nonzero higher-odd-inertia points: a finite
birational normalization map can have branches permuted by an
odd grading subgroup over its conductor locus.  Neither the
conductor nor the relevant coherent moving class has been computed
here.  The three-adic integrality of the entire normalization
series therefore does not remove $Y$ from the possible exception
in K-031 and does not prove a global pole at $512$.

\paragraph{Primary source and verification scope.}
Holweck, F., Luque, J.-G., and Thibon, J.-Y. (2012).
Geometric descriptions of entangled states by auxiliary varieties.
\emph{Journal of Mathematical Physics, 53}, 102203.
doi:10.1063/1.4753989; arXiv:1204.6375v2.
The complete binary-trilinear discussion in Section 5.2, printed
pp.~16--17, and the preceding character-series method in Section
5.1 have been read.  That source credits the classical covariant
system to Le Paige (1881); no original Le Paige proof is claimed
read.  Identity (K-033h) supplies a separate algebraic verification
of the exact series needed here.
